\nota{\textbf{Alcance documental.} Referencias transcritas del maestro y de su bibliografía local. El maestro declara corte normativo al 20/09/2026 y entrega al 21/09/2026. Esta edición no sustituye la verificación de fuentes originales: el Anexo B de vigencia citado por el maestro no está disponible en la carpeta recibida.}

\subsection{Fuentes legales y académicas}
\begin{enumerate}[label={[\arabic*]},leftmargin=2em,itemsep=4pt]
\item Congreso Nacional de Chile. (2024). \textit{Ley N.º 21.719: Regula la protección y el tratamiento de los datos personales y crea la Agencia de Protección de Datos Personales}. Diario Oficial, 13 de diciembre de 2024. \href{https://www.bcn.cl/leychile/navegar?i=1209272}{BCN, texto legal}.
\item Congreso Nacional de Chile. (2024). \textit{Ley N.º 21.663: Ley Marco de Ciberseguridad}. Diario Oficial, 8 de abril de 2024. \href{https://www.bcn.cl/leychile/navegar?idNorma=1202434}{BCN, texto legal}.
\item Congreso Nacional de Chile. (1999). \textit{Ley N.º 19.628: Sobre protección de la vida privada}. \href{https://www.bcn.cl/leychile/navegar?idNorma=141599}{BCN, texto legal}.
\item Ministerio del Interior y Seguridad Pública. (2024). \textit{Decreto Supremo N.º 295: Reglamento de reporte de incidentes de ciberseguridad}. Publicación indicada en el maestro: 1 de marzo de 2025. Canal operativo citado: \url{https://portal.anci.gob.cl}.
\item Parlamento Europeo y Consejo de la Unión Europea. (2016). \textit{Reglamento (UE) 2016/679 (RGPD)}. \href{https://eur-lex.europa.eu/legal-content/ES/TXT/?uri=CELEX:32016R0679}{EUR-Lex, CELEX 32016R0679}.
\item Parlamento Europeo y Consejo de la Unión Europea. (2022). \textit{Directiva (UE) 2022/2555 (NIS2)}. \href{https://eur-lex.europa.eu/legal-content/ES/TXT/?uri=CELEX:32022L2555}{EUR-Lex, CELEX 32022L2555}.
\item Parlamento Europeo y Consejo de la Unión Europea. (2024). \textit{Reglamento (UE) 2024/1689, de inteligencia artificial}. \href{https://eur-lex.europa.eu/legal-content/ES/TXT/?uri=CELEX:32024R1689}{EUR-Lex, CELEX 32024R1689}.
\item Parlamento Europeo y Consejo de la Unión Europea. (2024). \textit{Reglamento (UE) 2024/2847, de ciberresiliencia}. \href{https://eur-lex.europa.eu/legal-content/ES/TXT/?uri=CELEX:32024R2847}{EUR-Lex, CELEX 32024R2847}.
\item República Argentina. (2000). \textit{Ley N.º 25.326: Protección de los datos personales}. Fuente oficial indicada en el maestro: InfoLEG.
\item ISO/IEC. (2022). \textit{ISO/IEC 27001:2022. Sistemas de gestión de seguridad de la información}. \url{https://www.iso.org/standard/27001}.
\item ISO/IEC. (2025). \textit{ISO/IEC 27701:2025. Sistemas de gestión de privacidad}. \url{https://www.iso.org/standard/27701}.
\item ISO/IEC. (2023). \textit{ISO/IEC 42001:2023. Sistemas de gestión de inteligencia artificial}. \url{https://www.iso.org/standard/42001}.
\item NIST. (2024). \textit{The NIST Cybersecurity Framework (CSF) 2.0} (NIST CSWP 29). \url{https://doi.org/10.6028/NIST.CSWP.29}.
\item Escuela de Informática PUCV. (2026). \textit{Indicaciones Trabajo de Investigación 2026} (FEP00.3.26); \textit{Bases administrativas TFEP-01/2026} (E-24, E-25 y E-26); \textit{Bases técnicas del Caso 10: Transportes Curimón S.A.} (FEP03.10.26). Material del curso.
\end{enumerate}

\subsection{Registro del estado de vigencia consignado en el maestro}
\begin{tabularx}{\linewidth}{@{}p{3.6cm}Xp{2.8cm}@{}}
\toprule\textbf{Norma} & \textbf{Estado que informa la fuente} & \textbf{Corte declarado}\\\midrule
Ley 21.719 / Ley 19.628 & Reforma publicada; vigencia prevista 01/12/2026. Hasta entonces opera el régimen anterior. & 20/09/2026\\
Ley 21.663 / D.S. 295 & Operativos; reporte de incidentes y deberes de seguridad. & 20/09/2026\\
Ley 21.180 & Implementación gradual; tope legal señalado: 31/12/2027. & 20/09/2026\\
Marco europeo & Sin aplicación directa acreditada para Curimón; referencia de diseño y contratación. Calendario detallado pendiente del anexo de vigencia. & 20/09/2026\\
ISO y NIST & Ediciones y carácter certificable según página 5. & 20/09/2026\\\bottomrule
\end{tabularx}
